\chapter{Experiments}

\section{Experimental Setting}
This section summarizes the main components and configuration choices that define the experimental environment. It first presents the event logs used in the study and then outlines the embedding models, reranking component, and large language models considered in the evaluation.

To guide the experimental analysis, the study is organized around a small set of research questions framed in an exploratory way. In particular, the experiments examine: \textbf{RQ1}, how the proposed RAG-based approach behaves across different real-world event logs; \textbf{RQ2}, how prefix-construction choices such as \texttt{last-m}, \texttt{base}, and \texttt{gap} affect predictive performance; \textbf{RQ3}, how retrieval-side design choices, including embedding model and reranking, influence the final results; and \textbf{RQ4}, how the choice of LLM affects next-activity prediction quality under the proposed framework.

\subsection{Datasets}
The evaluation uses four publicly available real-life event logs that are widely used in predictive process monitoring and process mining research. Together, they cover financial, administrative, and healthcare processes, and they differ substantially in size, control-flow variability, and attribute richness. This diversity is useful for assessing whether the proposed retrieval-augmented approach remains effective across logs with different behavioral structure and operational context.

The \textit{BPI Challenge 2012} log captures the handling of personal loan and overdraft applications at a Dutch financial institution. It contains 262,200 events distributed across 13,087 traces. The process includes application registration, automated and manual checks, offer creation, offer acceptance, and final completion or cancellation, making it a strong benchmark for next-activity prediction under branching behavior and decision-intensive flows. In the local XES version used in this thesis, the log spans from October 1, 2011 to March 14, 2012, and it contains 24 distinct activity labels, or 36 activity classes when lifecycle transitions are taken into account \cite{bpi2012dataset}.

The \textit{BPI Challenge 2020 International Declarations} log contains travel-declaration and reimbursement cases related to international trips. It comprises 6,449 cases and 72,151 events. The process includes permit submission, approval steps by different organizational roles, declaration handling, and payment-related follow-up actions. Compared with the 2012 loan-application log, this dataset is smaller but still behaviorally rich, with repeated approval patterns, role-dependent routing, and multiple administrative checkpoints. In the local XES file used here, the observed event timestamps range from October 5, 2016 to May 10, 2020, and the log contains 34 distinct activity labels \cite{bpi2020intl}.

The \textit{Sepsis Cases} event log captures the trajectories of patients treated for sepsis in a hospital environment. It contains 1,050 traces and 15,214 events covering 16 distinct activities, along with a relatively rich set of recorded attributes, including organizational groups, laboratory-test results, and checklist-related information. Because the cases reflect real clinical pathways, the log is less regular than strongly structured administrative datasets and offers a useful benchmark for testing how the proposed method behaves under noisy and medically heterogeneous process executions. This intuition is also supported by recent evidence showing that the Sepsis Cases log exhibits increased complexity and variability, especially for longer prefixes, which makes next-activity prediction more challenging \cite{casciani2026rag}. In the local XES file used in this thesis, the timestamps range from November 7, 2013 to June 5, 2015 \cite{sepsisdataset}.

The \textit{Hospital Billing} event log was extracted from the financial modules of a regional hospital ERP system. It contains 100,000 traces and 451,359 events, recording the activities required to bill bundled medical-service packages. The log does not describe the clinical treatment itself; instead, it focuses on the administrative billing workflow. This makes it particularly useful for evaluating predictive methods in large-scale, repetitive back-office processes with moderate control-flow variation and substantial volume. In the present study, this dataset was not part of the baseline evaluation and was used mostly in the split-method evaluation. The local XES version spans from December 13, 2012 to January 19, 2016 and contains 18 distinct activity labels \cite{hospitalbillingdataset}.

\begin{table}[t]
\centering
\caption{Summary of the real-world event logs used in the evaluation}
\label{tab:dataset-summary}
\small
\setlength{\tabcolsep}{5pt}
\begin{tabular}{p{3.8cm}p{3.2cm}rrrr}
\toprule
Dataset & Time Span (months) & Traces & Events & Activity Labels \\
\midrule
BPI Challenge 2012 & 5.4 & 13{,}087 & 262{,}200 & 24 \\
BPI Challenge 2020\\(International Declarations) & 43.1 & 6{,}449 & 72{,}151 & 34 \\
Sepsis Cases & 18.9 & 1{,}050 & 15{,}214 & 16 \\
Hospital Billing & 37.2 & 100{,}000 & 451{,}359 & 18 \\
\bottomrule
\end{tabular}
\end{table}

\subsection{Embedding Models}
\begingroup
\emergencystretch=2em
In the retrieval component, two embedding models were examined. A MiniLM-based sentence-transformer model was selected as a reference point because it was also adopted in prior RAG-based work on next-activity prediction and yielded strong results there, making it an appropriate baseline for comparison \cite{casciani2026rag}. In parallel, a more recent and explicitly retrieval-oriented BGE model was evaluated in order to assess whether a newer embedding family designed more directly for dense retrieval could offer additional benefits. This comparison made it possible to contrast a strong prior reference encoder with a more recent retrieval-focused alternative under a similar retrieval pipeline.

The MiniLM configuration uses \texttt{sentence-transformers/all-MiniLM-L12-v2} \cite{minilmhf}. According to its model card, it maps sentences and short paragraphs to a 384-dimensional dense vector space and is intended for semantic search, clustering, and similarity-based retrieval, which makes it a reasonable lightweight baseline for prefix retrieval in process monitoring. Its underlying model family is rooted in MiniLM, which was designed through deep self-attention distillation to retain strong transformer behavior in a much smaller architecture \cite{minilmpaper}. In this thesis, MiniLM serves primarily as a compact baseline for testing whether a lighter encoder can still retrieve useful precedent prefixes.

By contrast, the BGE setting relies on \texttt{BAAI/bge-small-en-v1.5} \cite{bgesmallhf}. This model belongs to the BGE family, which is motivated by recent work on embedding models optimized for retrieval-oriented use cases, including strong multilinguality, multiple retrieval functions, and improved granularity through self-knowledge distillation \cite{bgem3paper}. Although the exact deployed model in the code is the smaller English \texttt{v1.5} variant, with a 384-dimensional representation space and a maximum input length of 512 tokens according to its model card, it follows the same broader design philosophy of the BGE line and is therefore a suitable retrieval-focused choice for the dense nearest-neighbor stage. In practice, it was selected as the main embedding model because it better reflects current retrieval-oriented embedding design while remaining lightweight enough for repeated experimentation.
\endgroup

\subsection{Reranker}
The reranking stage is a small extension of the retrieval pipeline. After the embedding model retrieves an initial pool of candidate prefixes from the vector database, a second model can be applied to reorder those candidates before the final top-$k$ examples are passed to the predictor. In the implementation used here, this stage relies on \texttt{BAAI/bge-reranker-base}, which is loaded as a sequence-classification cross-encoder and scores query--candidate pairs directly rather than comparing independent embeddings \cite{bgererankerhf}. This design follows the standard reranking setting introduced in neural passage reranking with BERT-style models \cite{nogueira2019bert}.

The main motivation for including this component was to test whether a more precise second-stage ranking step could improve the quality of the retrieved context with only a limited extension of the overall method. Since the reranker is applied only to a small candidate pool returned by dense retrieval, it preserves the general structure of the approach while offering a more targeted relevance assessment before prompt construction. In this way, reranking functions as a lightweight refinement of retrieval.

\subsection{Large Language Models}
The evaluation considered a mix of open-weight local models served through Ollama and proprietary models accessed through the OpenAI API. The goal was not to provide an exhaustive architectural analysis, but rather to compare models of different scale, design, and deployment setting under the same prediction task. Table~\ref{tab:llm-summary} summarizes the main LLMs used in the experiments together with the model information reported by their corresponding providers \cite{casciani2026rag,ollama_deepseek32b,ollama_gemma34b,ollama_gptoss20b,ollama_mistralsmall32_24b,ollama_phi414b,ollama_qwen330b,openai_gpt4omini,openai_gpt41mini}.

\begin{table}[H]
\centering
\caption{Summary of the LLMs considered in the evaluation}
\label{tab:llm-summary}
\small
\setlength{\tabcolsep}{5pt}
\begin{tabular*}{\textwidth}{@{\extracolsep{\fill}}p{3.9cm}p{1.5cm}p{2.6cm}p{1.8cm}p{1.8cm}p{1.6cm}}
\toprule
Model & Size & Architecture & Context & Embedding & Quant. \\
\midrule
\texttt{deepseek-r1:32b} & 32.8B & qwen2 & 131{,}072 & 5{,}120 & Q4\_K\_M \\
\texttt{gemma3:4b} & 4.3B & gemma3 & 131{,}072 & 2{,}560 & Q4\_K\_M \\
\texttt{gpt-oss:20b} & 20.9B & gptoss & 131{,}072 & 2{,}880 & MXFP4 \\
\texttt{mistral-small3.2:24b} & 24.0B & mistral3 & 131{,}072 & 5{,}120 & Q4\_K\_M \\
\texttt{phi4:14b} & 14.7B & phi3 & 16{,}384 & 5{,}120 & Q4\_K\_M \\
\texttt{qwen3:30b} & 30.5B & qwen3moe & 262{,}144 & 2{,}048 & Q4\_K\_M \\
\texttt{gpt-4o-mini} & N/P & Proprietary OpenAI & 128{,}000 & N/P & N/P \\
\texttt{gpt-4.1-mini-2025-04-14} & N/P & Proprietary OpenAI & 1{,}047{,}576 & N/P & N/P \\
\bottomrule
\end{tabular*}
\end{table}

The open-weight LLMs were served locally through Ollama (v0.30.10) on a workstation running Ubuntu 22.04.5 LTS. The platform featured an AMD Ryzen Threadripper PRO 5975WX processor with 32 cores, 251 GiB of system memory, and two NVIDIA GeForce RTX 4090 GPUs. Each GPU provided 24,564 MiB of VRAM, corresponding to approximately 24 GiB per device and about 48 GiB in total. These resources enabled the local execution of the quantized Ollama models used in the evaluation.

\section{Evaluation Metrics}
This section will describe accuracy, macro precision, macro recall, macro F1, and the earliness-oriented bucket analysis.

\section{Experimental Factors}
This section will summarize the controlled variables, such as top-k, gap, m, prompt variant, and reranking.

\section{Reproducibility}
This section will explain how runs, outputs, and result files are stored and compared.

\blankpage
